% !TEX program = xelatex
% Academic CV template for Jie Fu (傅杰)
% Compile with XeLaTeX.

\documentclass[11pt,a4paper]{article}

% ---------- Language and fonts ----------
\usepackage{iftex}
\ifXeTeX
  \usepackage{fontspec}
  \usepackage{xeCJK}
  \setmainfont{Times New Roman}
  \setsansfont{Arial}
  \setmonofont{Consolas}
  % If your TeX system cannot find this CJK font, replace it with another installed font.
  \setCJKmainfont{Noto Serif SC}
\else
  \usepackage[utf8]{inputenc}
  \usepackage[T1]{fontenc}
\fi

% ---------- Page setup ----------
\usepackage[a4paper,margin=1.9cm]{geometry}
\usepackage{enumitem}
\usepackage{tabularx}
\usepackage{xcolor}
\usepackage{hyperref}
\usepackage{titlesec}
\usepackage{array}
\usepackage{ragged2e}

\definecolor{linkblue}{RGB}{0,70,140}
\hypersetup{
  colorlinks=true,
  urlcolor=linkblue,
  linkcolor=linkblue,
  citecolor=linkblue,
  pdfauthor={Jie Fu},
  pdftitle={Curriculum Vitae -- Jie Fu}
}

\pagestyle{empty}
\setlength{\parindent}{0.5pt}
\setlength{\parskip}{2pt}
\setlist[itemize]{leftmargin=1.25em,itemsep=1pt,topsep=1.5pt}

\titleformat{\section}
  {\large\bfseries}
  {}{0pt}{}
  [\vspace{2pt}\titlerule]
\titlespacing*{\section}{0pt}{10pt}{6pt}

\newcolumntype{L}{>{\RaggedRight\arraybackslash}X}
\newcolumntype{R}{>{\RaggedLeft\arraybackslash}p{0.17\textwidth}}

\newcommand{\datedentry}[2]{%
  \begin{tabularx}{\textwidth}{R L}
    #1 & #2 \\
  \end{tabularx}\vspace{1pt}
}

\newcommand{\cvitem}[2]{%
  \begin{tabularx}{\textwidth}{>{\bfseries}p{0.35\textwidth} L}
    #1 & #2 \\
  \end{tabularx}
}

\begin{document}

% ---------- Header ----------
\begin{center}
  {\LARGE \textbf{Jie Fu} \; 傅杰}\\[3pt]
  Ph.D. Student in Mathematics, Qiuzhen College, Tsinghua University\\[3pt]
  \href{mailto:fu-j21@mails.tsinghua.edu.cn}{fu-j21@mails.tsinghua.edu.cn}
  \;\textbar\;
  \href{https://fujiemath.com}{fujiemath.com}
  % \;\textbar\;
  % \href{https://github.com/vegetablefj}{github.com/vegetablefj}
  \\[3pt]
  Tsinghua University, Beijing, China
\end{center}

\section{Personal Statement}

I am Jie Fu (傅杰 in Chinese).
%I was born in March 2006.
I am now a PhD student at Qiuzhen College, Tsinghua University.

% ---------- Research profile ----------
\section{Research Interests}
Algebraic geometry and related areas, with current work on automorphism groups of hypersurfaces and hyperkähler manifolds. Further interests include moduli theory, rationality problems and derived categories.

% ---------- Education ----------
\section{Education}
\datedentry{\textit{2025--present}}{%
\begin{tabular}{l}
     \textbf{Tsinghua University}, Qiuzhen College, Beijing, China.\\
      Ph.D. student in Mathematics.\\
      Advisor: Prof. Zhiwei Zheng.\\
\end{tabular}

}

\datedentry{\textit{2021--2025}}{%
  \begin{tabular}{l}
     \textbf{Tsinghua University}, Qiuzhen College, Beijing, China.\\
      Bachelor in Mathematics and Applied Mathematics.
\end{tabular}
}

% ---------- Publications and preprints ----------
\section{Publications and Preprints}
\subsection*{Preprints}
\begin{enumerate}[leftmargin=1.5em,label={[\arabic*]},itemsep=3pt]
  \item \textbf{Classification of Automorphism Groups of Smooth Cubic Threefolds and Fourfolds}.\\
  With Shihao Wang and Zhiwei Zheng. \emph{arXiv preprint}, arXiv:2609.20403 [math.AG], 2026. 86 pages.\\
  \href{https://arxiv.org/abs/2609.20403}{arxiv.org/abs/2609.20403}

  \item \textbf{Automorphism Groups of Smooth Cubic Fourfolds through Lattice Theory}.\\
  With Zhiwei Zheng. \emph{arXiv preprint}, arXiv:2609.06683 [math.AG], 2026. 26 pages.\\
  \href{https://arxiv.org/abs/2609.06683}{arxiv.org/abs/2609.06683}

  \item \textbf{Non-symplectic Indices of Automorphism Groups of Smooth Cubic Fourfolds}.\\
  With Shihao Wang and Zhiwei Zheng. \emph{arXiv preprint}, arXiv:2606.11754 [math.AG], 2026. 30 pages.\\
  \href{https://arxiv.org/abs/2606.11754}{arxiv.org/abs/2606.11754}
\end{enumerate}

% \subsection*{Journal Publications}
% No peer-reviewed journal publications yet.

% ---------- Selected notes ----------
% \section{Selected Notes and Expository Writing}
% % TODO: Keep this section only if you want lecture notes/expository notes on your CV.
% \begin{itemize}
%   \item \emph{Notes on Finiteness, Linearity and Generic Triviality of Automorphism Groups of Smooth Hypersurfaces}.
%   \item \emph{Notes on Addictive Decomposition of Polynomials}. % TODO: Check whether ``Addictive'' should be ``Additive''.
%   \item \emph{Hypersurface with Automorphism Group of Large Order}.
% \end{itemize}

% ---------- Talks ----------
% \section{Talks}
% \datedentry{20XX}{\textbf{TODO: Talk title}. TODO: Seminar/Conference, Institution, City.}
% \datedentry{20XX}{\textbf{TODO: Talk title}. TODO: Seminar/Conference, Institution, City.}

% ---------- Academic activities ----------
\newpage
\section{Talks, Seminars and Academic Activities}

\begin{itemize}
    \item \textit{Smooth hypersurface with automorphism group of large order}. Spring Learning Seminar in Algebraic Geometry, YMSC, Beijing. March 18th, 2026.
    \item \textit{New results on automorphism groups of cubic fourfolds}. YMSC Seminar on Moduli Spaces and Related Topics, YMSC, Beijing. March 12th, 2026.
    \item \textit{Symplectic automorphism of $K3^{[n]}$-type manifolds}. Hyperkähler Seminar, YMSC, Beijing. December 2nd, 2025.
    \item \textit{Linear system of K3 surfaces}. Hyperkähler Seminar, YMSC, Beijing. October 24th, 2025.
    \item \textit{Basic lattice theory following Nikulin}. Hyperkähler Seminar, YMSC, Beijing. September 30th, 2025.
    \item Co-organization: Student Seminar on Perverse Sheaves and Application to Representation Theory, Qiuzhen College, Beijing, Fall 2025.
\end{itemize}

% ---------- Teaching ----------
\section{Teaching Experience}
\datedentry{Spring 2026}{ Teaching assistant, \textit{Algebraic Variety} given by Prof. Junpeng Jiao, Tsinghua University.}

% ---------- Honors ----------
 \section{Honors and Awards}
 \datedentry{Spring 2026}{Pi Day's best problems award, Qiuzhen College.}
 \datedentry{Spring 2023}{Honorable Mention, Analysis and Differential Equations, S.-T. Yau College Student Mathematics Contest.}
  \datedentry{Fall 2020}{First prize in China's national high school mathematics competition, Chinese Mathematical Society.}
% \datedentry{20XX}{TODO: Award name, granting institution.}

% ---------- Service ----------
% \section{Academic Service}
% \datedentry{20XX}{TODO: Refereeing, seminar organization, outreach, or departmental service.}

% ---------- Skills ----------
% \section{Technical Skills}
% %\cvitem{Languages}{TODO: Chinese; English; other languages.}

% %\cvitem{Basic Programming}{\TeX; C/C++; Python; HTML; R.}

% \cvitem{Computational Software}{GAP; SageMath; Magma; OSCAR/Julia.}


% ---------- References ----------
\section{References}
\begin{tabularx}{\textwidth}{L l}
\textbf{Prof. Zhiwei Zheng} & Advisor.\\
 & Assistant Professor,\\
 & Tsinghua University, Yau Mathematical Sciences Center, \\
 & zhengzhiwei@mail.tsinghua.edu.cn
\end{tabularx}

\vfill
\begin{center}
  {\small Last updated: \today}
\end{center}

\end{document}
